\uuid{r6cP}
\exo7id{1852}
\titre{exo7 1852}
\auteur{maillot}
\organisation{exo7}
\datecreate{2001-09-01}
\isIndication{false}
\isCorrection{false}
\chapitre{Equation aux dérivées partielles}
\sousChapitre{Equation aux dérivées partielles}
\cpge{mpsi-mp2i.raisonnement-vocabulaire-ensembliste.c, pcsi.raisonnement-vocabulaire-ensembliste.d, ptsi.raisonnement-vocabulaire-ensembliste.d, mpsi-mp2i.derivabilite.d, pcsi.derivabilite.d, ptsi.derivabilite.d, mpsi-mp2i.fonctions-deux-variables.b, pcsi.fonctions-deux-variables.b, ptsi.fonctions-deux-variables.b, mpsi-mp2i.fonctions-deux-variables.c, pcsi.fonctions-deux-variables.c, ptsi.fonctions-deux-variables.c, bcpst1.vocabulaire-applications, bcpst1.derivees-primitives-integrales.b, bcpst1.fonctions-deux-variables.c, mp-mpi.calcul-differentiel-optimisation.a}
\module{Equation différentielle}
\niveau{L3}
\difficulte{}

\contenu{
\texte{
On cherche les fonctions $f:\R^2\rightarrow \R$ telles que:
\begin{equation}\label{truc}
\frac{\partial f}{\partial u} (u,v) +2u \frac{\partial f}{\partial v} (u,v)=0 \qquad \text{pour tout $(u,v)\in\R^2$}.
\end{equation}
Soit $\phi:\R^2\rightarrow\R^2$ l'application d\'efinie par $\phi(x,y)=(x,y+x^2)$.
}
\begin{enumerate}
    \item \question{En calculant l'application r\'eciproque, montrer que $\phi$ est
bijective. V\'erifier que $\phi$ et $\phi^{-1}$ sont de classe $\mathcal C^1$.}
    \item \question{Soit $f:\R^2\rightarrow \R$ une fonction de classe $\mathcal C^1$. Posons
$g=f\circ \phi$.
  }
    \item \question{Montrer que $g$ est de classe $\mathcal C^1$.}
    \item \question{Montrer que $f$ est solution de (\ref{truc}) si et seulement si
  $\frac{\partial g}{\partial x}=0$.}
\end{enumerate}
}
